\onehalfspacing
\phantomsection
\addcontentsline{toc}{section}{\textit{Abstract}}
{\fontsize{16pt}{30pt}\selectfont \noindent\textbf{Abstract}}
\vspace{0.5cm}
{\fontsize{12pt}{18pt}\selectfont

Epilepsy is a prevalent neurological disorder affecting approximately 50 million people worldwide. Automated analysis of electroencephalogram (EEG) signals may help reduce the burden of manual seizure review. This study proposes a hybrid deep learning architecture combining a 1D Convolutional Neural Network (1D-CNN), Bidirectional Long Short-Term Memory network (Bi-LSTM), and Self-Attention mechanism for automated epileptic seizure detection using the University of Bonn EEG dataset. A strict recording-level 10-Fold cross-validation strategy is employed to prevent data leakage from overlapping sliding-window segments. The proposed model is systematically compared against three baselines --- SVM with wavelet-packet features, a pure 1D-CNN, and a pure Bi-LSTM with attention --- in an ablation study across binary, three-class, and five-class classification tasks. Experimental results show that the hybrid model achieves 99.67\% accuracy on binary classification, 97.69\% on three-class, and 84.51\% on the five-class task. On the five-class task, it shows statistically significant gains over the pure CNN ($p = 0.001$) and pure Bi-LSTM ($p = 0.001$) baselines, while the mean gain over SVM does not reach significance at the 0.05 level. Attention weight visualisation highlights salient temporal regions and provides a qualitative view of model behaviour.

}
\vspace{0.5cm}
\addcontentsline{toc}{section}{\textit{Keywords}}
{\fontsize{16pt}{30pt}\selectfont \noindent\textbf{Keywords}}
\vspace{0.5cm}
{\fontsize{12pt}{18pt}\selectfont

Epilepsy Detection, Electroencephalogram (EEG), Deep Learning, 1D Convolutional Neural Network, Bidirectional LSTM, Self-Attention, Bonn Dataset, Cross-Validation

}
